\section{Only the base tuple is needed for the complex Bessel identity}

% BEGIN R89C1
\paragraph{R89C1. The selected spectral model is available at every sufficiently late index.}
Fix a family in the selected class of X87S1, in the small positive
$\epsilon$ range of P88A1. There is a single threshold such that the
spectral model of its actual selected bundle satisfies the full
condition $\mathrm{FinX}$ at every subsequent index. In particular its
comparison of the abstract period space with its automorphic realization
can be evaluated there.

Indeed TFR2, independently of the volume estimate in that proposition,
gives $n\ge1801$ and gaps of at least $16n^{19}$ between successive
bottom weights, at every place and every index beyond one threshold.
These are the two conditions used in SFS4. The supplied local Bessel
wave is reproduced by the discrete-series projector. Jiang--Sun--Zhu,
\emph{Uniqueness of Bessel models: the Archimedean case}, Theorem A,
gives local multiplicity at most one for the generic Bessel character
and compact target identified in SFS3. Thus the product projector
has image of dimension at most one on Bessel sections, and the normalized wave has
the same absolute value as the explicit decaying wave for that model.
SFS4 gives the weighted integrability and reproduction for all words
and integrability of the common kernel. SFS5 then applies to the
spectral space $X$ and the arithmetic data $Y$ of the selected
bundle itself. In particular its period series belong to the specified
Hilbert space, its column maps are linear, and its evaluations,
fibres, Gram parameters and measures have the required compatibilities.
Taking the maximum with the selection threshold and
$\mathcal F.x_1+1$ gives the assertion for every $i$ at once.
This disposes of the spectral admissibility input in the period
comparison when a family in the class is supplied.
% END R89C1

% BEGIN R89C2
\paragraph{R89C2. A base-tuple identity suffices for all complex arguments.}
Use the selected model of R89C1. Suppose its arithmetic and spectral
realizations have the original comparison of row coordinates, column
coordinates and inner products ($\mathrm{JoinAt}$). Let $t_0$ be the
base tuple, $k\hookrightarrow\mathbb C$ the model embedding $\tau$,
and $m(u)$ the column label of word $u$. Write
\[
  Z_j=\Theta_{t_0} j_R(e_{V,t_0} e_j),\qquad
  p_{u,ij}=\tau\bigl(\mathrm{Pt}_{m(u),t_0}(e_i,e_j)\bigr),
\]
using the coordinate vectors in the two indicated rational spaces.
It suffices to produce one unitary $A$ on the good model row space
and, for each word, a unitary $N_u$ in its model column metric
and a scalar $a_u\ne0$ such that
\[
  p_{u,ij}=a_u\,
  \big\langle s_{EE}(A Z_j),P_{m(u)}(N_u e_i)\big\rangle
                  \qquad\text{for all }i,j.                 \tag{C89.1}
\]
Here $P_{m(u)}$ is the model period operator and $s_{EE}$ is its
conjugate-linear row realization. In this implication the $N_u$ and
$a_u$ may be chosen word by word even if two words have the same model
label. No tuple other than $t_0$ occurs in (C89.1).

To prove the implication, let $J$ be the row coordinate map from the
native row space onto the good model row space. The join makes it a
positive scalar times a semilinear isometry. Its semilinearity is
given by $\eta$, either identity or complex conjugation. The column
coordinate map $K_u$ followed by the same $\eta$ applied to each
coordinate is likewise an isometry up to a positive scalar into the
column metric. Set
\[
  V=J^{-1}AJ,\qquad
  U_u=(\eta K_u)^{-1}N_u(\eta K_u).
\]
The two occurrences of $\eta$ cancel and the positive scales divide
out, so these are complex-linear unitaries. The same $V$ works for
every word. For $c_F$ the native base-tuple input coordinate map, the
row comparison gives $Z_j=\lambda J\iota(c_F e_j)$ with one
$\lambda\ne0$. The chain comparison identifies the native arithmetic
evaluation on the basis pair with $\eta(p_{u,ij})$. Inserting the row
comparison in (C89.1) thus contributes the scalar $\eta(a_u\lambda)$, since
$s_{EE}$ and the first slot of the inner product are both
conjugate-linear. The inner-product comparison for this word contributes
its fixed nonzero factor $\kappa_u$ in passing to the native spectral
pairing. All coordinate vectors $e_i,e_j$ are fixed by $\eta$.
Both resulting native expressions are complex bilinear in their
coordinate vectors, so equality on these pairs gives equality on the
whole product.

Finally the join's input coordinates cover the good native input space.
By R88L1--2 its image under $\iota$ is the full row target, and
$\iota$ restricts to an isomorphism on the orthogonal complement of its
kernel there. For $F$ in the row target take this distinguished lift.
Its arithmetic evaluation is $B_u(\phi,F)$ by the definition of $B_u$.
The spectral column comparison may compose $U_u$ with a further
wordwise unitary. This proves exactly
$B_u(\phi,F)=\eta(a_u\lambda)\kappa_u\,
 \mathrm{besS}_u(U'_u\phi,VF)$, the complex clause of entry four.
The proof of the implication uses neither an assertion on the
non-base factors of the mixed identity nor an estimate on the sizes
of its nonzero scalars. The stronger all-tuple comparison in the source
also serves arithmetic purposes elsewhere.
% END R89C2

% BEGIN R89C3
\paragraph{R89C3. The relative periods can be tested by rational vanishing.}
For proposed $A,N_u$ in R89C2 fix a word and let $s_{ij}$ be the inner
product in (C89.1), without the scalar. The arithmetic array $p$ has
entries in $k_0=\tau(k)$. If $p$ is identically zero, the test is
exactly that $s$ vanish. Otherwise choose any index pair $b$ with
$p_b\ne0$. There exists $a_u\ne0$ as in (C89.1) if and only if
$s_b\ne0$ and the spectral evaluation vanishes on the following
explicit rational relations for all other index pairs $l$, namely
\[
                 p_b\,s_l-p_l\,s_b=0.                       \tag{C89.2}
\]
Indeed the tensors $p_b e_l-p_l e_b$ span over $k_0$ the kernel of
the arithmetic functional on the coordinate tensor product. To see
this, write any vector in this kernel in the coordinate basis and
use its one linear equation to eliminate the coefficient at $b$.
These same vectors span its complex kernel after extending scalars.
If (C89.2) holds, $a_u=p_b/s_b$ works at all entries. The converse
follows by substitution. A zero-dimensional coordinate space is covered
by the identically zero case.

Consequently the residual comparison has a formulation suitable for a
theta identity. With one row isometry for all words, it suffices to show that
the analytic integrals kill the explicit arithmetic relations, with a
nonzero value at one arithmetic pivot for each nonzero word. The
rationality of the arithmetic coefficients is what makes the relations
available over the coefficient field. Verifying these vanishings and the pivot values for the selected data
remains part of the global period comparison.
% END R89C3
